\documentclass[11pt]{article}
\usepackage{fontspec}
\setmainfont{Times New Roman}
\setsansfont{Arial}
\setmonofont{Consolas}
\usepackage{amsmath,amssymb,mathtools}
\usepackage{geometry}
\usepackage{microtype}
\newcommand{\mo}{\mathfrak{o}}
\newcommand{\mO}{\mathfrak{O}}
\newcommand{\mT}{\mathfrak{T}}
\newcommand{\mR}{\mathfrak{R}}
\newcommand{\mq}{\mathfrak{q}}
\geometry{a4paper,margin=2.35cm}
\setlength{\parindent}{1.25em}
\setlength{\parskip}{0pt}
\makeatletter
\renewcommand\footnotesize{\@setfontsize\footnotesize{7.4}{8.4}\selectfont}
\makeatother
\emergencystretch=100em
\allowdisplaybreaks
\sloppy
\hbadness=10000
\vbadness=10000
\hfuzz=2pt
\vfuzz=10pt
\raggedbottom

\begin{document}

% Unit: Paper32 second introductory paragraph v001. Direct Russian translation from clean RA34 Paper32 line 10 / segment P32-S0005, checked against source scan pages 1--2 and Claude batch OCR witness.
\setcounter{footnote}{0}

\noindent В дальнейшем вопрос о полях расщепления будет прослежен дальше. Поля расщепления наименьшей степени можно характеризовать посредством соответствующих некоммутативных тел; общие поля расщепления -- посредством двусторонне простых колец, инвариантно связанных с этими некоммутативными телами. Далее на примере тела кватернионов будет показано, что степени минимальных полей расщепления не ограничены; при этом поле расщепления называется минимальным, если ни одно из его собственных подполей не является полем расщепления. Эта неограниченность существенно следует из теоретико-числовой теоремы существования, доказательство которой дано Г. Хассе в непосредственно следующей заметке. Однако пример также разбирается элементарно-вычислительно -- скорее в проверочном плане; тем самым без привлечения общей теории и теоретико-числовой теоремы существования неограниченность степеней показывается путем элементарного доказательства менее сильной, но достаточной теоремы существования.\footnote{Этот пример принадлежит Э. Нётер; он основан на модификации примера Р. Брауэра, в котором вообще было показано существование минимального поля расщепления, степень которого превышает индекс. Элементарная обработка примера принадлежит Р. Брауэру. Характеризация полей расщепления восходит к исследованиям авторов, которые в остальном преследуют раздельные цели. У Э. Нётер речь идет о построении теории представлений на основе теории модулей и идеалов при общих предположениях конечности; у Р. Брауэра -- о построении некоммутативных тел конечной степени над заданным коммутативным совершенным основным полем с помощью фактор-систем. К характеризации полей расщепления наименьшей степени авторы пришли независимо: Р. Брауэр для совершенного основного поля, Э. Нётер без этого ограничения. Характеризацию общих полей расщепления Р. Брауэр нашел в связи с вопросом Э. Нётер о возможности некоммутативного вложения; Э. Нётер и здесь смогла снять ограничение на совершенную основную область.} Следует еще заметить, что и в этом примере речь идет о полях расщепления конечной группы, а именно группы кватернионов, лежащей также в основе примера Шура. Указанные там представления являются одновременно (изоморфными) представлениями тела кватернионов.

\end{document}
